\section{Neural Rendering 背景}

在讨论 Score Distillation 之前，需要理解一个关键前提——\textbf{可微分渲染}（differentiable rendering）。

\subsection{从多视图图像到 3D 重建}

Neural Rendering 的基本任务是：给定一个场景的多视角图像集合，重建出该场景的 3D 表示，使得从任意新视角渲染出的图像都尽可能逼真。

\begin{knowledgebox}{Neural Rendering 的核心思想}
Neural Rendering 使用\textbf{神经网络}作为 3D 场景的表示。网络参数 $\theta$ 隐式地编码了场景的几何与外观信息。给定一个相机视角参数 $\pi$，通过可微分的渲染函数 $g(\theta, \pi)$ 可以生成对应的 2D 图像。

代表性工作包括 NeRF（Neural Radiance Fields）和 3D Gaussian Splatting。
\end{knowledgebox}

\subsection{传统 3D 重建的优化目标}

给定 $N$ 张已知视角的参考图像 $\{I_i\}_{i=1}^{N}$，3D 重建的优化目标非常直接：

$$
\theta^* = \arg\min_{\theta} \sum_{i=1}^{N} \| g(\theta, \pi_i) - I_i \|_2^2
$$

其中：
\begin{itemize}
    \item $\theta$：3D 表示的参数（如 NeRF 的网络权重）
    \item $\pi_i$：第 $i$ 张图像的相机参数（位姿）
    \item $g(\theta, \pi_i)$：从视角 $\pi_i$ 渲染得到的图像
    \item $I_i$：对应视角的真实参考图像
\end{itemize}

\begin{warningbox}{重建质量取决于输入图像数量}
传统的 Neural Rendering 需要较多的输入图像（通常 $\sim$100 张）才能获得高质量的 3D 重建。当图像数量减少到 10 张、5 张甚至 1 张时，重建质量会急剧下降——未覆盖视角的区域会出现明显的伪影（如黑色区域、模糊等）。
\end{warningbox}

\subsection{实际应用中的挑战}

Neural Rendering 在实际应用中面临诸多挑战：

\begin{itemize}
    \item \textbf{重建速度}：计算量大，不适合实时应用
    \item \textbf{渲染速度}：需要快速渲染以支持交互式浏览
    \item \textbf{场景规模}：从单个物体到城市级场景
    \item \textbf{动态物体}：如何处理场景中的运动物体
    \item \textbf{输入图像数量}：如何用尽可能少的图像实现高质量重建
\end{itemize}

特别是在电商场景中（如在线购物的 3D 产品展示），用户期望仅拍摄少量照片就能快速生成 3D 模型，这要求极少输入图像下的高质量重建。

\subsection{本章小结}

Neural Rendering 提供了一个从 2D 图像到 3D 表示的桥梁，其核心是\textbf{可微分渲染函数}。传统方法依赖大量输入图像，而引入外部先验（如预训练扩散模型的知识）可以显著减少所需图像数量，甚至实现零样本 3D 生成。

\newpage
